\section{Conclusion}

Ce chantier a commenc\'e par une anomalie : deux chemins qui partagent le
m\^eme r\'eseau, le m\^eme GPU et le m\^eme \'evaluateur, mais un facteur
quatre de vitesse par position. Il se termine par une correction de quelques
lignes, s\'emantiquement neutre, qui ram\`ene la recherche UCI au niveau du
self-play.

La d\'emarche compte autant que le r\'esultat. Un banc hors arbre a mesur\'e
l'\'evaluateur seul, avec un mode qui reproduit fid\`element la fa\c con dont
la recherche alterne les tailles de lot. Ce mode a r\'ev\'el\'e que
\code{session->Run} passe de 3 \`a 13\,ms selon que la forme du lot est stable
ou non. Le reste en d\'ecoule : padder les lots \`a une forme fixe, ignorer les
sorties dupliqu\'ees, et laisser la racine au lot 1.

Les mesures interleaved donnent $\times 2{,}0$ en ouverture, $\times 3{,}2$ en
milieu et $\times 2{,}55$ en finale pour huit workers, avec un p95 de latence
divis\'e par deux \`a trois. Le m\^eme gain s'applique au chemin mono, donc au
banc de puzzles et aux tournois. Le bot UCI active le r\'eglage.

Deux pistes ont \'et\'e refus\'ees sur preuves : FP16, plus lent de 18 \`a
28\,\%, et les r\'eglages de divergence, plus rapides de 13 \`a 25\,\% mais
moins pr\'ecis de 1,8 \`a 2,4 points. Cette seconde r\'efutation est aussi
instructive que le gain : elle montre que la barri\`ere qualit\'e du projet
fonctionne et qu'un d\'ebit plus \'elev\'e n'est pas une fin en soi.

Enfin, la mesure du cas chaud a corrig\'e une erreur de raisonnement des
rapports pr\'ec\'edents : un succ\`es de table n'\'evite pas l'inf\'erence, car
la descente continue vers une feuille r\'eseau. La r\'eutilisation d'arbre
approfondit la recherche, elle n'all\`ege pas le travail GPU, et la r\'eserve
Amdahl est lev\'ee.

Il reste du chemin vers Leela Chess Zero : une file d'inf\'erence continue
pour remplir des lots plus gros ($\times 1{,}3$ \`a $\times 1{,}7$), et un
backend mieux adapt\'e aux cartes 8$\times$8. Le prochain jalon naturel est un
microbenchmark de file, puis un plan s\'epar\'e pour la production continue.
